% !TEX root = ./main.tex

% setup.tex
\fdusetup{
  style = {
    font-size = -4,
    bib-backend = biblatex,        
    bib-resource = {main.bib},     
    bib-style = gb7714-2025        % 更新参考文献样式
  },
  info = {
    title = {论文题目},
    title* = {English Title},
    author = {某\quad 某},
    author* = {Jane Doe},
    supervisor = {某\quad 某\quad 教~授},
    department = {数学科学学院},
    major = {计算数学},
    student-id = {10010101010},
    instructors = {{某\quad 某\quad 教~授},{某\quad 某\quad 教~授},{某某某\quad 教~授}},
    keywords = {关键词1, 关键词2, 关键词3, 关键词4},
    keywords* = {KeyWords1, KeyWords2, KeyWords3, KeyWords4},
    clc = {O24}
  }
}
% 字体，英文字体更新为Times New Roman
\setmainfont{Times New Roman}
% 序号
\usepackage{enumerate}

% 定义图片文件目录与扩展名
\graphicspath{{figs/}}
\DeclareGraphicsExtensions{.pdf,.eps,.png,.jpg,.jpeg}

% 子图宏包
\usepackage{subcaption}
\usepackage{bicaption}

% 确定浮动对象的位置，可以使用 [H]，强制将浮动对象放到这里（可能效果很差）
\usepackage{float}

% 表格排列
\usepackage{array}

% 固定宽度的表格
\usepackage{tabularx}

% 使用三线表：toprule，midrule，bottomrule。
\usepackage{booktabs}

% 表格中支持跨行
\usepackage{multirow}

% 使用长表格
\usepackage{longtable}

% 附带脚注的表格
\usepackage{threeparttable}

% 附带脚注的长表格
\usepackage{threeparttablex}

% 算法环境宏包
\usepackage[ruled,vlined,linesnumbered]{algorithm2e}
% \usepackage{algorithm, algorithmicx, algpseudocode}

% 代码环境宏包
\usepackage{listings}
\lstdefinestyle{lstStyleCode}{%
  aboveskip         = \medskipamount,
  belowskip         = \medskipamount,
  basicstyle        = \ttfamily\zihao{6},
  commentstyle      = \slshape\color{black!60},
  stringstyle       = \color{green!40!black!100},
  keywordstyle      = \bfseries\color{blue!50!black},
  extendedchars     = false,
  upquote           = true,
  tabsize           = 2,
  showstringspaces  = false,
  xleftmargin       = 1em,
  xrightmargin      = 1em,
  breaklines        = false,
  framexleftmargin  = 1em,
  framexrightmargin = 1em,
  columns           = flexible,
  keepspaces        = true,
  texcl             = true,
  mathescape        = true
}

% 绘图宏包
\usepackage{tikz}
\usetikzlibrary{arrows.meta, shapes.geometric,cd}

% 超链接
\usepackage{hyperref}
%更新区------------------------------------------------------------------------
\lstnewenvironment{codeblock}[1][]{%
  \lstset{style=lstStyleCode,#1}}{}
%封面字体更新
\AtBeginDocument{
 \addtocontents{toc}{\protect\hypersetup{linkcolor=black}}
  \addtocontents{toc}{\protect\color{black}}

  \addtocontents{lof}{\protect\hypersetup{linkcolor=black}}
  \addtocontents{lof}{\protect\color{black}}

  \addtocontents{lot}{\protect\hypersetup{linkcolor=black}}
  \addtocontents{lot}{\protect\color{black}}

  \addtocontents{loa}{\protect\hypersetup{linkcolor=black}}
  \addtocontents{loa}{\protect\color{black}}

  \addtocontents{lol}{\protect\hypersetup{linkcolor=black}}
  \addtocontents{lol}{\protect\color{black}}
}
%摘要更新
\ExplSyntaxOn
\cs_set_protected:Npn \__fdu_abstract_begin:
  {
    \cleardoublepage
    \thispagestyle { plain }
    % —— 补回被丢掉的目录项 + PDF 书签 ——
    \phantomsection
    \addcontentsline { toc } { chapter } { \c__fdu_name_abstract_tl }
    % —— 补回页眉 ——
    \__fdu_chapter_header:n { \c__fdu_name_abstract_tl }
    % —— 你的自定义排版 ——
    \zihao { -4 }
    \begin { center }
      \l__fdu_info_title_tl
    \end { center }
    \noindent \c__fdu_name_abstract_tl \c__fdu_fwid_colon_tl
  }
\cs_set_protected:Npn \__fdu_abstract_en_begin:
  {
    \cleardoublepage
    \thispagestyle { plain }
    \phantomsection
    \addcontentsline { toc } { chapter } { \c__fdu_name_abstract_en_tl }
    \__fdu_chapter_header:n { \c__fdu_name_abstract_en_tl }
    \zihao { -4 }
    \begin { center }
      \l__fdu_info_title_en_tl
    \end { center }
    \noindent \c__fdu_name_abstract_en_tl :
  }
\ExplSyntaxOff

\ExplSyntaxOn
\cs_set_protected:Npn \__fdu_abstract_end:
  {
    \__fdu_keywords:nNn
      { \sffamily \c__fdu_name_keywords_tl \c__fdu_fwid_colon_tl }
      \l__fdu_info_keywords_clist { \c__fdu_fwid_semicolon_tl }
  }
\ExplSyntaxOff

\ExplSyntaxOn
\cs_set_protected:Npn \__fdu_abstract_en_end:
  {
    \__fdu_keywords:nNn
      { \bfseries \c__fdu_name_keywords_en_tl \__fdu_quad: }
      \l__fdu_info_keywords_en_clist { ; ~ }
  }
\ExplSyntaxOff
%封面更新
\ExplSyntaxOn
\cs_set_protected:Npn \__fdu_cover_id:
  {
    \noindent
    \begin{minipage}[t]{0.5\textwidth}
      \raggedright
      中图分类号：\l__fdu_info_clc_tl
    \end{minipage}%
    \begin{minipage}[t]{0.5\textwidth}
      \raggedleft
      学校代码：10246 \\
      % 同等学力人员/英文项目   % ← 需要时去掉本行行首的 % 即可显示
    \end{minipage}
  }
\ExplSyntaxOff

\ExplSyntaxOn
\cs_set_protected:Npn \__fdu_cover_info:
  {
    \begin{center}
    \zihao{-3}        % ← 封面字号
    \begin{tabular}{ @{} r @{：} c @{} }
      \makebox[7em][s]{学生姓名} & \l__fdu_info_author_tl       \\[1ex]
      \makebox[7em][s]{学号}     & \l__fdu_info_student_id_tl   \\[1ex]
      \makebox[7em][s]{导师姓名、职称} & \l__fdu_info_supervisor_tl \\[1ex]
      \makebox[7em][s]{学科专业} & \l__fdu_info_major_tl        \\[1ex]
      \makebox[7em][s]{培养单位} & \l__fdu_info_department_tl   \\[1ex]
      \makebox[7em][s]{论文提交日期} & \quad \the\year 年 \the\month 月 \the\day 日 \\
    \end{tabular}
    \end{center}
  }
\ExplSyntaxOff